% generated by GAPDoc2LaTeX from XML source (Frank Luebeck)
\documentclass[a4paper,11pt]{report}

\usepackage[top=37mm,bottom=37mm,left=27mm,right=27mm]{geometry}
\sloppy
\pagestyle{myheadings}
\usepackage{amssymb}
\usepackage[utf8]{inputenc}
\usepackage{makeidx}
\makeindex
\usepackage{color}
\definecolor{FireBrick}{rgb}{0.5812,0.0074,0.0083}
\definecolor{RoyalBlue}{rgb}{0.0236,0.0894,0.6179}
\definecolor{RoyalGreen}{rgb}{0.0236,0.6179,0.0894}
\definecolor{RoyalRed}{rgb}{0.6179,0.0236,0.0894}
\definecolor{LightBlue}{rgb}{0.8544,0.9511,1.0000}
\definecolor{Black}{rgb}{0.0,0.0,0.0}

\definecolor{linkColor}{rgb}{0.0,0.0,0.554}
\definecolor{citeColor}{rgb}{0.0,0.0,0.554}
\definecolor{fileColor}{rgb}{0.0,0.0,0.554}
\definecolor{urlColor}{rgb}{0.0,0.0,0.554}
\definecolor{promptColor}{rgb}{0.0,0.0,0.589}
\definecolor{brkpromptColor}{rgb}{0.589,0.0,0.0}
\definecolor{gapinputColor}{rgb}{0.589,0.0,0.0}
\definecolor{gapoutputColor}{rgb}{0.0,0.0,0.0}

%%  for a long time these were red and blue by default,
%%  now black, but keep variables to overwrite
\definecolor{FuncColor}{rgb}{0.0,0.0,0.0}
%% strange name because of pdflatex bug:
\definecolor{Chapter }{rgb}{0.0,0.0,0.0}
\definecolor{DarkOlive}{rgb}{0.1047,0.2412,0.0064}


\usepackage{fancyvrb}

\usepackage{mathptmx,helvet}
\usepackage[T1]{fontenc}
\usepackage{textcomp}


\usepackage[
            pdftex=true,
            bookmarks=true,        
            a4paper=true,
            pdftitle={Written with GAPDoc},
            pdfcreator={LaTeX with hyperref package / GAPDoc},
            colorlinks=true,
            backref=page,
            breaklinks=true,
            linkcolor=linkColor,
            citecolor=citeColor,
            filecolor=fileColor,
            urlcolor=urlColor,
            pdfpagemode={UseNone}, 
           ]{hyperref}

\newcommand{\maintitlesize}{\fontsize{50}{55}\selectfont}

% write page numbers to a .pnr log file for online help
\newwrite\pagenrlog
\immediate\openout\pagenrlog =\jobname.pnr
\immediate\write\pagenrlog{PAGENRS := [}
\newcommand{\logpage}[1]{\protect\write\pagenrlog{#1, \thepage,}}
%% were never documented, give conflicts with some additional packages

\newcommand{\GAP}{\textsf{GAP}}

%% nicer description environments, allows long labels
\usepackage{enumitem}
\setdescription{style=nextline}

%% depth of toc
\setcounter{tocdepth}{1}





%% command for ColorPrompt style examples
\newcommand{\gapprompt}[1]{\color{promptColor}{\bfseries #1}}
\newcommand{\gapbrkprompt}[1]{\color{brkpromptColor}{\bfseries #1}}
\newcommand{\gapinput}[1]{\color{gapinputColor}{#1}}


\begin{document}

\logpage{[ 0, 0, 0 ]}
\begin{titlepage}
\mbox{}\vfill

\begin{center}{\maintitlesize \textbf{ PackageMaker \mbox{}}}\\
\vfill

\hypersetup{pdftitle= PackageMaker }
\markright{\scriptsize \mbox{}\hfill  PackageMaker  \hfill\mbox{}}
{\Huge \textbf{ A \textsf{GAP} package for creating new \textsf{GAP} packages \mbox{}}}\\
\vfill

{\Huge  1.0.0 \mbox{}}\\[1cm]
{ 24 March 2026 \mbox{}}\\[1cm]
\mbox{}\\[2cm]
{\Large \textbf{ Max Horn\\
    \mbox{}}}\\
\hypersetup{pdfauthor= Max Horn\\
    }
\end{center}\vfill

\mbox{}\\
{\mbox{}\\
\small \noindent \textbf{ Max Horn\\
    }  Email: \href{mailto://mhorn@rptu.de} {\texttt{mhorn@rptu.de}}\\
  Homepage: \href{https://www.quendi.de/math} {\texttt{https://www.quendi.de/math}}\\
  Address: \begin{minipage}[t]{8cm}\noindent
 Fachbereich Mathematik\\
 RPTU Kaiserslautern\texttt{\symbol{45}}Landau\\
 Gottlieb\texttt{\symbol{45}}Daimler\texttt{\symbol{45}}Stra{\ss}e 48\\
 67663 Kaiserslautern\\
 Germany\\
 \end{minipage}
}\\
\end{titlepage}

\newpage\setcounter{page}{2}
\newpage

\def\contentsname{Contents\logpage{[ 0, 0, 1 ]}}

\tableofcontents
\newpage

     
\chapter{\textcolor{Chapter }{PackageMaker}}\label{Chapter_PackageMaker}
\logpage{[ 1, 0, 0 ]}
\hyperdef{L}{X80AC704D873187E3}{}
{
  
\section{\textcolor{Chapter }{What is \textsf{PackageMaker}?}}\label{Chapter_PackageMaker_Section_What_is_PackageMaker}
\logpage{[ 1, 1, 0 ]}
\hyperdef{L}{X78E9A6927A4E1FBD}{}
{
  \textsf{PackageMaker} is a \textsf{GAP} package that makes it easy to create new \textsf{GAP} packages, by providing a ``wizard'' function that asks a few questions about the new package, such as its name and
authors; and from that creates a complete usable skeleton package. It
optionally can set up several advanced features, including: 
\begin{itemize}
\item  a git repository for use on GitHub; 
\item  a simple C or C++ kernel extension, including a build system; 
\item  continuous integration and release automation using GitHub Actions; 
\item  code coverage tracking; 
\item  ... and more. 
\end{itemize}
 

 The generated package is meant as a starting point, not as a finished package
ready for release. In particular, you should expect to edit \texttt{README.md}, \texttt{PackageInfo.g}, the generated tests and, depending on your choices, the documentation, CI
workflows or kernel extension sources. }

 
\section{\textcolor{Chapter }{Before you start}}\label{Chapter_PackageMaker_Section_Before_you_start}
\logpage{[ 1, 2, 0 ]}
\hyperdef{L}{X7E730E5A7DD77A50}{}
{
  A typical session looks like this: 
\begin{itemize}
\item  start \textsf{GAP} in a directory where it is OK to create a new package directory; 
\item  load \textsf{PackageMaker} and run \texttt{PackageWizard()}; 
\item  answer the wizard questions; 
\item  inspect and edit the generated files; 
\item  move the package to a directory in which \textsf{GAP} looks for packages, or run \textsf{GAP} with \texttt{\texttt{\symbol{45}}\texttt{\symbol{45}}packagedirs}. 
\end{itemize}
 

 By default, the wizard creates the new package directory in the current
working directory of your \textsf{GAP} session. That is often convenient for creating the package, but it does not
automatically make the package loadable. See the next section for the
different ways to make \textsf{GAP} find the generated package. 

 If you ask \textsf{PackageMaker} to create a git repository, git should have \texttt{user.name} and \texttt{user.email} configured. If they are missing, \textsf{PackageMaker} explains how to set them up, lets you retry, or keeps the generated package
directory without creating a git repository. 

 }

 
\section{\textcolor{Chapter }{Where should the generated package go?}}\label{Section_PackageDestination}
\logpage{[ 1, 3, 0 ]}
\hyperdef{L}{X7BB8C09C87932665}{}
{
  The generated package directory must eventually live in a place where \textsf{GAP} searches for packages. 

 For normal day\texttt{\symbol{45}}to\texttt{\symbol{45}}day use, the best
destination is usually \texttt{\texttt{\symbol{126}}/.gap/pkg}. You may have to create that directory first. Keeping personal packages there
means they survive \textsf{GAP} upgrades and are separate from the main \textsf{GAP} installation. 

 Another natural choice is the \texttt{pkg} directory inside your main \textsf{GAP} installation. This is common when you built \textsf{GAP} from source in a directory that you control. The drawback is that when you
later upgrade or reinstall \textsf{GAP}, it is easy to forget to copy over packages that you added there manually. 

 For quick testing, \textsf{GAP} 4.15 or later can be started with \texttt{gap \texttt{\symbol{45}}\texttt{\symbol{45}}packagedirs .} to tell \textsf{GAP} to search the current directory directly for packages. 

 There are also more advanced alternatives: 
\begin{itemize}
\item  command line options such as \texttt{\texttt{\symbol{45}}l} (see  (\textbf{Reference: Command Line Options})) can change the configured \textsf{GAP} root directories (see  (\textbf{Reference: GAP Root Directories})) when starting \textsf{GAP}; 
\item  \texttt{ExtendRootDirectories} (\textbf{Reference: ExtendRootDirectories}) and \texttt{ExtendPackageDirectories} (\textbf{Reference: ExtendPackageDirectories}) can add package search locations after \textsf{GAP} has already started; 
\item  \texttt{SetPackagePath} (\textbf{Reference: SetPackagePath}) can be useful when you want to force loading one particular package from a
specific path for a session. 
\end{itemize}
 

 These advanced options are useful to know about, but for most users the
practical advice is simple: keep the package in \texttt{\texttt{\symbol{126}}/.gap/pkg} for normal use, and use \texttt{gap \texttt{\symbol{45}}\texttt{\symbol{45}}packagedirs .} only as a quick\texttt{\symbol{45}}testing convenience when running \textsf{GAP} 4.15 or newer. }

 
\section{\textcolor{Chapter }{A worked example}}\label{Chapter_PackageMaker_Section_A_worked_example}
\logpage{[ 1, 4, 0 ]}
\hyperdef{L}{X858C9C7F832B9C2C}{}
{
  The following transcript shows a complete example session. It uses GitHub,
enables the generated workflows, chooses the MIT license, creates a C kernel
extension, and enters one author / maintainer record. 

 

\subsection{\textcolor{Chapter }{PackageWizardInput}}
\logpage{[ 1, 4, 1 ]}\nobreak
\hyperdef{L}{X7ECA3E578388AD9A}{}
{\noindent\textcolor{FuncColor}{$\triangleright$\enspace\texttt{PackageWizardInput({\mdseries\slshape arg})\index{PackageWizardInput@\texttt{PackageWizardInput}}
\label{PackageWizardInput}
}\hfill{\scriptsize (function)}}\\


 Interactively create a package skeleton. You can abort by either answering ``quit'' or pressing \texttt{Ctrl\texttt{\symbol{45}}D}. 

 
\begin{Verbatim}[commandchars=!|B,fontsize=\small,frame=single,label=Example]
  !gapprompt|gap>B !gapinput|PackageWizard();B
  Welcome to the GAP PackageMaker Wizard.
  I will now guide you step-by-step through the package
  creation process by asking you some questions.
  
  What is the name of the package? DemoPackage
  Enter a short (one sentence) description of your package: A package used to demonstrate PackageMaker
  Which license should the package use? MIT
  Shall I prepare your new package for GitHub? [Y/n] y
  Do you want to use GitHub Actions for automated tests and making releases? [Y/n] y
  The release workflow updates the package website on GitHub Pages
  whenever you make a package release.
  I need to know the URL of the GitHub repository.
  It is of the form https://github.com/USER/REPOS.
  What is USER (typically your GitHub username)? demo-user
  What is REPOS, the repository name? DemoPackage
  Shall your package provide a GAP kernel extension? Yes, written in C
  
  Next I will ask you about the package authors and maintainers.
  
  Last name? Doe
  First name(s)? Dana
  Is this one of the package authors? [Y/n] y
  Is this a package maintainer? [Y/n] y
  Email? dana@example.invalid
  WWWHome? https://example.invalid/~dana/
  GitHubUsername? demo-user
  PostalAddress? Example Institute\nDepartment of Algebra\nExample Street 1\n12345 Example City
  Place? Example City, Country
  Institution? Example Institute
  Add another person? [y/N] n
  Creating the git repository...
  Initialized empty Git repository in .../DemoPackage/.git/
  [main (root-commit) ...] initial import
  Done creating git repository.
  Create <https://github.com/demo-user/DemoPackage> via <https://github.com/new> and then run:
    git push -u origin main
\end{Verbatim}
 

 After this finishes, the new package lives in a directory named \texttt{DemoPackage} below the current directory. You can now inspect and edit the generated files,
move the package into a package directory, and try loading it in \textsf{GAP}. }

 

\subsection{\textcolor{Chapter }{PackageWizardGenerate}}
\logpage{[ 1, 4, 2 ]}\nobreak
\hyperdef{L}{X86E254658700A465}{}
{\noindent\textcolor{FuncColor}{$\triangleright$\enspace\texttt{PackageWizardGenerate({\mdseries\slshape arg})\index{PackageWizardGenerate@\texttt{PackageWizardGenerate}}
\label{PackageWizardGenerate}
}\hfill{\scriptsize (function)}}\\


 

 }

 

\subsection{\textcolor{Chapter }{PackageWizard}}
\logpage{[ 1, 4, 3 ]}\nobreak
\hyperdef{L}{X7B19742884A7E49C}{}
{\noindent\textcolor{FuncColor}{$\triangleright$\enspace\texttt{PackageWizard({\mdseries\slshape arg})\index{PackageWizard@\texttt{PackageWizard}}
\label{PackageWizard}
}\hfill{\scriptsize (function)}}\\


 

 }

 }

 
\section{\textcolor{Chapter }{Important wizard choices}}\label{Chapter_PackageMaker_Section_Important_wizard_choices}
\logpage{[ 1, 5, 0 ]}
\hyperdef{L}{X87B127887E7D3B43}{}
{
  Some answers mainly affect metadata, while others determine which files are
generated. 

 
\subsection{\textcolor{Chapter }{Package name}}\label{Chapter_PackageMaker_Section_Important_wizard_choices_Subsection_Package_name}
\logpage{[ 1, 5, 1 ]}
\hyperdef{L}{X7B45A0267CB58132}{}
{
  The package name becomes the name of the top\texttt{\symbol{45}}level package
directory and is also used in file names such as \texttt{gap/YourPackage.gd}, \texttt{gap/YourPackage.gi} and in many entries of \texttt{PackageInfo.g}. It therefore should be chosen with care. 

 }

 
\subsection{\textcolor{Chapter }{License}}\label{Chapter_PackageMaker_Section_Important_wizard_choices_Subsection_License}
\logpage{[ 1, 5, 2 ]}
\hyperdef{L}{X861E5DF986F89AE2}{}
{
  The license choice affects both the generated \texttt{LICENSE} file and the guidance inserted into \texttt{README.md}. By default, the wizard uses \texttt{GPL\texttt{\symbol{45}}2.0\texttt{\symbol{45}}or\texttt{\symbol{45}}later}, because that is the license of \textsf{GAP} itself and is also used by many \textsf{GAP} packages. Other built\texttt{\symbol{45}}in choices are \texttt{GPL\texttt{\symbol{45}}3.0\texttt{\symbol{45}}or\texttt{\symbol{45}}later}, \texttt{MIT} and \texttt{BSD\texttt{\symbol{45}}3\texttt{\symbol{45}}Clause}. 

 If you are unsure which license to use, then the following sites may help: 
\begin{itemize}
\item  \href{https://choosealicense.com/} {\texttt{https://choosealicense.com/}} 
\item  \href{https://opensource.org/licenses/} {\texttt{https://opensource.org/licenses/}} 
\item  \href{https://spdx.org/licenses/} {\texttt{https://spdx.org/licenses/}} 
\end{itemize}
 

 If you choose the custom option, then \textsf{PackageMaker} generates a placeholder \texttt{LICENSE} file which you must replace with the full text of your chosen license. You
should also make sure that the \texttt{License} field in \texttt{PackageInfo.g} and the package \texttt{README.md} describe the same license. 

 }

 
\subsection{\textcolor{Chapter }{GiHub setup}}\label{Chapter_PackageMaker_Section_Important_wizard_choices_Subsection_GiHub_setup}
\logpage{[ 1, 5, 3 ]}
\hyperdef{L}{X812040BA822BCC79}{}
{
  Saying yes to GitHub setup does two things: 
\begin{itemize}
\item  it fills in repository\texttt{\symbol{45}}related URLs in \texttt{PackageInfo.g}; 
\item  it offers to create a local git repository with an initial commit and a GitHub
remote named \texttt{origin}. 
\end{itemize}
 

 If you also enable GitHub Actions, then \textsf{PackageMaker} adds CI and release workflows below \texttt{.github/workflows/} and creates \texttt{.codecov.yml}. The release workflow updates the package website on GitHub Pages when you
make a release. 

 If you do not use GitHub setup, then \texttt{PackageInfo.g} contains placeholder URL entries instead. They work as reminders, but you
should replace them before publishing your package. 

 }

 
\subsection{\textcolor{Chapter }{Kernel extension}}\label{Chapter_PackageMaker_Section_Important_wizard_choices_Subsection_Kernel_extension}
\logpage{[ 1, 5, 4 ]}
\hyperdef{L}{X857DC0D182FFA23F}{}
{
  The kernel extension choice determines whether the package is purely
interpreted \textsf{GAP} code or also contains code in \texttt{C} or \texttt{C++}. Choosing a kernel extension adds files such as \texttt{src/YourPackage.c} or \texttt{src/YourPackage.cc}, together with \texttt{configure}, \texttt{Makefile.in} and \texttt{Makefile.gappkg}. }

 }

 
\section{\textcolor{Chapter }{What gets generated?}}\label{Chapter_PackageMaker_Section_What_gets_generated}
\logpage{[ 1, 6, 0 ]}
\hyperdef{L}{X7D674D9B85B3F6C3}{}
{
  Every generated package contains the basic files needed for a small but usable
package: 
\begin{itemize}
\item  \texttt{PackageInfo.g} with the package metadata; 
\item  \texttt{README.md} with placeholder text to replace; 
\item  \texttt{init.g} and \texttt{read.g}; 
\item  \texttt{gap/YourPackage.gd} and \texttt{gap/YourPackage.gi}; 
\item  \texttt{makedoc.g} for building the package manual; 
\item  \texttt{tst/testall.g} as a starting point for tests; 
\item  \texttt{LICENSE}. 
\end{itemize}
 

 Depending on your choices, the wizard may also create: 
\begin{itemize}
\item  \texttt{.gitattributes} and \texttt{.gitignore} when GitHub setup is enabled; 
\item  \texttt{.github/workflows/CI.yml}, \texttt{.github/workflows/release.yml} and \texttt{.codecov.yml} when GitHub Actions are enabled; 
\item  \texttt{src/}, \texttt{configure}, \texttt{Makefile.in} and \texttt{Makefile.gappkg} when a kernel extension is requested. 
\end{itemize}
 

 A few of these files deserve immediate attention: 
\begin{itemize}
\item  \texttt{README.md} still contains TODO text and should be rewritten; 
\item  \texttt{PackageInfo.g} should be checked carefully, especially the package URLs, description, authors
and maintainers; 
\item  \texttt{tst/testall.g} is only the test driver, so you will usually want to add actual \texttt{.tst} files; 
\item  the generated \texttt{gap/} and optional \texttt{src/} files are only a starting skeleton and should be adapted to your package. 
\end{itemize}
 }

 
\section{\textcolor{Chapter }{What next?}}\label{Chapter_PackageMaker_Section_What_next}
\logpage{[ 1, 7, 0 ]}
\hyperdef{L}{X868D9A7B7BAAE5D3}{}
{
  Once the wizard has finished, a typical next\texttt{\symbol{45}}step checklist
is: 
\begin{itemize}
\item  decide where the package should live, as explained in \ref{Section_PackageDestination}; 
\item  edit \texttt{README.md} and replace the TODO sections; 
\item  review \texttt{PackageInfo.g} and fix any remaining placeholder values; 
\item  if you selected the custom license option, replace the placeholder \texttt{LICENSE} text with the full license text; 
\item  add package code and tests; 
\item  run the package inside \textsf{GAP} and make sure it loads; 
\item  if you asked for a kernel extension, run \texttt{./configure} and build it; 
\item  if you enabled GitHub setup, create the GitHub repository and push the
generated initial commit. 
\end{itemize}
 

 For broader guidance on package structure and release preparation, it is worth
reading the \textsf{GAP} manual chapter on  (\textbf{Reference: Using and Developing GAP Packages}) as well as the manual and \texttt{PackageInfo.g} file of the \href{https://github.com/gap-packages/example} {\texttt{https://github.com/gap\texttt{\symbol{45}}packages/example}} package. }

 }

 \def\indexname{Index\logpage{[ "Ind", 0, 0 ]}
\hyperdef{L}{X83A0356F839C696F}{}
}

\cleardoublepage
\phantomsection
\addcontentsline{toc}{chapter}{Index}


\printindex

\immediate\write\pagenrlog{["Ind", 0, 0], \arabic{page},}
\newpage
\immediate\write\pagenrlog{["End"], \arabic{page}];}
\immediate\closeout\pagenrlog
\end{document}
